% ================= 第 4 章 (后半): 4.3 + 4.4 =================

\subsection{(N,p) = 1 时定理 1.1.3 的证明}\label{sec:4.3}

保留 \ref{sec:4.2} 的记号, 这里再假设 $(N,p) = 1$.
固定一个如 \ref{sec:3.3} 的模形式 $f$, 使得 $\sigma_p(f)$ 特别是晶体的.
在这种情况下, $\chi_p$ 是把 $p$ 送到 $\chi(p)^{-1}$ 的 $\Qp^{\times}$ 的非分歧特征标, $\chi_2$ 非分歧, 并设 $\alpha_p := \chi_2(p)$, $\beta_p := p^{k-1}\chi_1(p)$.
有 $\val(\beta_p) = 0$, $\alpha_p\beta_p = \chi(p)p^{k-1}$, 且由 \ref{sec:3.3}:
\[
\begin{aligned}
&D_p(f) = Le_1 \oplus Le_2,\\
&\phi(e_1) = \alpha_p e_1,\\
&\phi(e_2) = \beta_p e_2,\\
&\mathrm{Fil}^{i}D_p(f) = D_p(f), \qquad i \le 0,\\
&\mathrm{Fil}^{i}D_p(f) = L(e_1 + \delta e_2), \qquad 1 \le i \le k-1,\\
&\mathrm{Fil}^{i}D_p(f) = 0, \qquad i \ge k,
\end{aligned}
\]
其中 $\delta = 0$ 当且仅当 $\sigma_p(f)$ 分裂.

先假设 $N > 4$.
滤过 $\phi$-模 $D_p(f)$ 此时有如下几何描述 (参见 e.g.\ \cite[\S1]{ref:29}):
\[
D_p(f) \simeq H^{1}\bigl(X_1(N),H_{k-2} \otimes \Omega^{\bullet}(\log)\bigr) \otimes_{\Qp} L,
\qquad
\mathrm{Fil}^{i}D_p(f) \simeq H^{0}(X_1(N),\omega^{k}) \otimes_{\Qp} L \quad (1 \le i \le k-1),
\]
即 $L \cdot f$ (记住 $f$ 是新形式), 其中 $D_p(f)$ 上的 Frobenius $\phi$ 来自 $H^{1}\bigl(X_1(N),H_{k-2} \otimes \Omega^{\bullet}(\log)\bigr)$ 上的晶体 Frobenius $\phi$.

设 $\mathcal{X}$ 为 $X_1(N)$ 在 $\Zp$ 上的紧光滑模型, $A \in H^{0}(\mathcal{X} \times_{\Zp} \mathbb{F}_p,\omega^{p-1})$ 为 Hasse 不变量.
回顾 $A$ 在 $\mathcal{X} \times_{\Zp} \mathbb{F}_p$ 的每个超奇异点以 $1$ 阶零点消失 (且在其它点不消失).
若 $x \in X_1(N)(\Qp)$, 记 $\bar{x}$ 为 $x$ 在 $(\mathcal{X} \times_{\Zp} \mathbb{F}_p)(\mathbb{F}_p)$ 中的特殊化.
如 \cite{ref:10} 附录中所解释, 对每个 $\eta \in p^{\Q}$ 且 $p^{-1} \le \eta < 1$, 可以定义 $X_1(N)$ 的一个刚性解析拟 Stein 开集 $W(\eta)$:
\[
x \in W(\eta) \iff |A(\bar{x})| > \eta.
\]
若 $p > 3$, 则 $W(\eta) \subset X_1(N)$ 是 $X_1(N)$ 的满足下式的闭点 $x$ 组成的刚性解析开集:
\[
x \in W(\eta) \iff |E_{p-1}(x)| > \eta,
\]
其中 $E_{p-1} \in H^{0}(X_1(N),\omega^{p-1})$ 是权 $p-1 \ge 4$、水平 $1$ 的 Eisenstein 级数 (参见 \cite[\S1 和 \S2]{ref:13}).

设 $W_1 := W\bigl(p^{-p/(p+1)}\bigr)$.
由于 $W_1$ 拟 Stein, 有:
\[
H^{1}\bigl(W_1,H_{k-2} \otimes \Omega^{\bullet}(\log)\bigr) = \frac{H^{0}\bigl(W_1,H_{k-2} \otimes \Omega^{1}(\log C)\bigr)}{\nabla H^{0}(W_1,H_{k-2})},
\]
且有交换图:
\begin{equation}
\label{eq:18}
\begin{CD}
H^{0}(W_1,\omega^{k}) @>>> H^{1}\bigl(W_1,H_{k-2} \otimes \Omega^{\bullet}(\log)\bigr)\\
@AAA @AAA\\
H^{0}(X_1(N),\omega^{k}) @>{ }>> H^{1}\bigl(X_1(N),H_{k-2} \otimes \Omega^{\bullet}(\log)\bigr)
\end{CD}
\end{equation}
其中两个竖直箭头都是单射 (对第二个, 容易由 \cite[Th.~2.1 和 Th.~2.4]{ref:2} 推出).
设 $W_2 := W\bigl(p^{-1/(p+1)}\bigr) \subset W_1$.
记 $E_2$ 和 $E_1$ 分别为 $W_2$ 和 $W_1$ 之上的泛广义椭圆曲线, 则有一个交换图 (参见 \cite[\S2]{ref:13}):
\begin{equation}
\label{eq:19}
\begin{CD}
E_2 @>{\Phi}>> E_1\\
@VVV @VVV\\
W_2 @>{\varphi}>> W_1
\end{CD}
\end{equation}
它由函子性诱导"Frobenius"映射:
\[
\varphi : H^{0}(W_1,\omega^{k-2} \otimes \Omega^{1}(\log C)) \to H^{0}(W_2,\omega^{k-2} \otimes \Omega^{1}(\log C)),
\]
\[
\Phi : H^{1}\bigl(W_1,H_{k-2} \otimes \Omega^{\bullet}(\log)\bigr) \to H^{1}\bigl(W_2,H_{k-2} \otimes \Omega^{\bullet}(\log)\bigr).
\]
从"模空间"的观点看, 映射 $\varphi : W_2 \to W_1$ 把 $(E,P)$ (其中 $E$ 是 $W_2$ 的一点之上的广义椭圆曲线, $P$ 是 $E$ 上精确 $N$ 阶的有理点) 送到 $(E/\mathrm{can},P)$, 其中 $\mathrm{can} \subset E$ 是 $E$ 的标准子群 (\cite[\S3]{ref:35}), $P$ 是 $P$ 在商 $E/\mathrm{can}$ 中的像.

\begin{lemma}\label{lem:4.3.1}
下图:
\[
\begin{CD}
H^{1}\bigl(X_1(N),H_{k-2}\otimes\Omega^{\bullet}(\log)\bigr) @>{ }>> H^{1}\bigl(W_1,H_{k-2}\otimes\Omega^{\bullet}(\log)\bigr) @<<< H^{0}(W_1,\omega^{k-2}\otimes\Omega^{1}(\log C))\\
@V{\phi}VV @VV{\Phi}V @VV{\varphi}V\\
H^{1}\bigl(X_1(N),H_{k-2}\otimes\Omega^{\bullet}(\log)\bigr) @>{ }>> H^{1}\bigl(W_2,H_{k-2}\otimes\Omega^{\bullet}(\log)\bigr) @<<< H^{0}(W_2,\omega^{k-2}\otimes\Omega^{1}(\log C))
\end{CD}
\]
是交换的.
\end{lemma}

\begin{proof}
右边方块的交换是显然的.
给出左边方块交换的主要论证.
为简化记号, 记 $X = X_1(N)$, $\mathcal{X}$ 为 $X$ 在 $\Zp$ 上的紧光滑模型, $\widehat{\mathcal{X}}$ 为其 $p$ 进形式完备化, $\bar{X}$ 为它们在 $\mathbb{F}_p$ 上的特殊纤维, $Z \subset \bar{X}$ 为序点开集, $\mathcal{Z} \subset \mathcal{X}$ 为 $Z$ 在 $\mathcal{X}$ 中按 \cite[\S1]{ref:5} 意义的管 $]Z[_{\widehat{\mathcal{X}}}$.
它是 $\mathcal{X}$ 的一个仿射型集.
考虑 $\widehat{\mathcal{X}}$ 之上的泛 (广义) 椭圆曲线, 可见紧簇 $X$ 上的凝聚层 $H_{k-2}$ 来自 $\widehat{\mathcal{X}}$ 上的凝聚层, 即晶体 $E_{k-2}$ 在加厚 $X \hookrightarrow \widehat{\mathcal{X}}$ 上的求值, 其中 $E_{k-2}$ 是晶体位形 $(\bar{X}/\mathbb{F}_p)_{\mathrm{cris}}$ 上的晶体 (更精确地, 需要考虑尖点处的对数结构, 并把 $E_{k-2}$ 定义为 $\bar{X}/\mathbb{F}_p$ 的对数晶体位形上的对数晶体, 对数结构在 $\mathbb{F}_p$ 上平凡; 但为简化行文, 我们在此证明中处处忽略这一点, 它不会带来问题).
群 $H^{1}\bigl(X_1(N),H_{k-2} \otimes \Omega^{\bullet}(\log)\bigr)$ 正是 $H^{1}_{\mathrm{cris}}(\bar{X},E_{k-2}) \otimes \Qp$.
事实上, $E_{k-2}$ 是 \cite[\S2.4.1]{ref:5} 意义下的非退化 $F$-晶体 (利用 $X$ 之上泛椭圆曲线上的 Frobenius), 并由限制诱导沿 $X - Z$ 的超收敛 $F$-等晶体, 记为 $E^{\dagger}_{k-2}$ (参见 \cite[\S2.3]{ref:5}).
对满足 $p^{-1/(p+1)} < \eta \le 1$ 的 $\eta \in p^{\Q}$, 设 $Z(\eta) \subset W_2$ 为 $\mathcal{X}$ 的仿射型集, 其闭点 $x$ 满足 $|A(\bar{x})| \ge \eta$.
有 $Z(1) = Z$, 且 $Z(\eta)_{\eta<1}$ 构成 $Z$ 在 $\mathcal{X}$ 中按 \cite[\S1.2]{ref:5} 意义的严格邻域的基本系.
此外, 对 $\eta > p^{-1/p(p+1)}$ 有 $\varphi(Z(\eta)) \subset Z(\eta^{p}) \subset W_2$.
由 \cite[Th.~2.4]{ref:2}, 对 $\eta \le \eta'$, 限制映射 $H^{1}\bigl(Z(\eta),H_{k-2} \otimes \Omega^{\bullet}(\log)\bigr) \to H^{1}\bigl(Z(\eta'),H_{k-2} \otimes \Omega^{\bullet}(\log)\bigr)$ 是同构.
设:
\[
H^{1}_{\mathrm{rig}}\bigl(Z,E^{\dagger}_{k-2}\bigr) := \varinjlim_{\eta<1} H^{1}\bigl(Z(\eta),H_{k-2} \otimes \Omega^{\bullet}(\log)\bigr)
\]
配备其 Frobenius 自同态 $\Phi$, 由 Frobenius $\varphi : Z(\eta) \to Z(\eta^{p})$ (以及泛椭圆曲线上的 Frobenius) 诱导.
左边方块的交换即等价于证明限制映射:
\[
H^{1}_{\mathrm{cris}}(\bar{X},E_{k-2}) \otimes \Qp \to H^{1}_{\mathrm{rig}}\bigl(Z,E^{\dagger}_{k-2}\bigr)
\]
与 Frobenius 交换, 但这是显然的: 在晶体一侧和刚性一侧, Frobenius 自同态都通过函子性得到, 只要借助 $\widehat{\mathcal{X}}$ 的一个带局部 Frobenius 提升的开覆盖上的双重 de Rham 复形计算两个上同调即可.
\end{proof}

回顾: $H^{0}(W_1,\omega^{k})$ 的任何截面 $g$ 在 $W_1$ 的每个尖点 $c$ 处有 $q$-展开 $\sum_{n=0}^{\infty} a_n(g)(c)q^{n} \in \Qp \otimes \Zp[[q]]$, 由 $g$ 在 $W_1$ 在尖点 $c$ 处的形式完备化之上的泛广义椭圆曲线处求值给出.

\begin{lemma}\label{lem:4.3.2}
设 $g \in H^{0}\bigl(X_1(N),\omega^{k-2} \otimes \Omega^{1}(\log C)\bigr) \simeq H^{0}(X_1(N),\omega^{k})$, 记 $\phi(g) \in H^{1}\bigl(X_1(N),H_{k-2} \otimes \Omega^{\bullet}(\log)\bigr)$ 为 $g$ 经晶体 Frobenius 的像.
则 $\phi(g)$ 经图 \eqref{eq:18} 来自 $H^{0}(W_1,\omega^{k})$ 中一个 $q$-展开在无穷远尖点处为:
\[
\sum_{n=0}^{+\infty} a_n(\phi(g))(\infty)q^{n} = p^{k-1} \sum_{n=0}^{+\infty} a_n(\langle p \rangle g)(\infty)q^{pn}
\]
的截面.
特别地, 若 $\langle p \rangle g = \chi(p)g$, 则 $\phi(g)$ 是 $p^{k-1}\chi(p)g(pz)|_{W_1}$ 的像.
\end{lemma}

\begin{proof}
由引理 \ref{lem:4.3.1}, 只需验证: 若 $g$ 是 $H^{0}\bigl(W_1,\omega^{k-2} \otimes \Omega^{1}(\log C)\bigr)$ 中的形式, 记 $g(\infty)(q)$ 为 $g$ 在无穷远尖点处的 $q$-展开, 则 $\varphi(g) \in H^{0}\bigl(W_2,\omega^{k-2} \otimes \Omega^{1}(\log C)\bigr)$ 满足:
\[
\varphi(g)(\infty)(q) = p^{k-1}g(\langle p \rangle \infty)(q^{p})
\]
(因为 $\sum_{n=0}^{+\infty} a_n(\langle p \rangle g)(\infty)q^{pn} = \sum_{n=0}^{+\infty} a_n(g)(\langle p \rangle \infty)q^{pn}$).
为此, 在无穷远尖点处的完备化之上把截面 $\varphi(g)$ 求值在 Tate 曲线 $(\mathbb{G}_m/q^{\mathbb{Z}},\zeta_N)$ (带其 $N$ 阶点) 上, 用 \cite[Prop.~3.1.3(1)]{ref:29} 的 (标准) 记号, 并利用 $\mathrm{can}(\mathbb{G}_m/q^{\mathbb{Z}}) = \langle \zeta_p \rangle$, 得到:
\[
\varphi(g)(\infty)(q) \otimes \frac{dq}{q}\Bigl(\frac{dt}{t}\Bigr)^{k-2}
= \varphi(g)\bigl(\mathbb{G}_m/q^{\mathbb{Z}},\zeta_N\bigr)(q) \otimes \frac{dq}{q}\Bigl(\frac{dt}{t}\Bigr)^{k-2}
\]
\[
= g\bigl(\mathbb{G}_m/\langle \zeta_p,q^{\mathbb{Z}} \rangle,\zeta_N\bigr)(q) \otimes \frac{dq}{q}\Bigl(\frac{dt}{t}\Bigr)^{k-2}
= g\bigl(\mathbb{G}_m/q^{p\mathbb{Z}},\zeta_N^{p}\bigr)(q^{p}) \otimes \frac{dq^{p}}{q^{p}}\Bigl(\frac{dt}{t^{p}}\Bigr)^{k-2}
\]
\[
= p^{k-1}g(\langle p \rangle \infty)(q^{p}) \otimes \frac{dq}{q}\Bigl(\frac{dt}{t}\Bigr)^{k-2},
\]
由此得到结果.
\end{proof}

回顾 (参见 \ref{sec:4.2}) 有 $\theta^{k-1} : H^{0}(W_1,\omega^{2-k}) \to H^{0}(W_1,\omega^{k})$, 且 $f(z) - \beta_p f(pz)$ 水平为 $\Gamma_1(Np)$、是 $U_p$ 关于特征值 $\alpha_p$ 的特征向量 (参见 \ref{sec:4.1}).

\begin{theorem}\label{thm:4.3.3}
设 $f$ 是对 $\Gamma_1(N)$ 为新的权 $k \ge 2$ 本征尖点形式, $(N,p) = 1$, 且局部表示 $\sigma_p(f)$ 可约.
若 $N > 4$, 则 $\sigma_p(f)$ 分裂当且仅当存在 $g \in H^{0}(W_1,\omega^{2-k}) \otimes L$ 使得在 $H^{0}(W_1,\omega^{k}) \otimes L$ 中 $f(z) - \beta_p f(pz) = \theta^{k-1}(g)$.
若 $N \le 4$, 取 $q \nmid Mp$ 为足够大的素数, 则 (用 $X_1(N;q)$ 代替 $X_1(N)$, 参见 \ref{sec:4.2}) $\sigma_p(f)$ 分裂当且仅当存在 $g \in H^{0}(W_1,\omega^{2-k})^{\GL_2(\mathbb{Z}/q\mathbb{Z})} \otimes_{\Qp} L$ 使得在 $H^{0}(W_1,\omega^{k})^{\GL_2(\mathbb{Z}/q\mathbb{Z})} \otimes_{\Qp} L$ 中 $f(z) - \beta_p f(pz) = \theta^{k-1}(g)$.
\end{theorem}

\begin{proof}
表示 $\sigma_p(f)$ 分裂当且仅当在 $D_p(f)$ 的前述描述中, 有 $\mathrm{Fil}^{k-1}D_p(f) = Le_1$ 且 $\phi(e_1) = \alpha_p e_1$.
先假设 $N > 4$.
在 $D_p(f)$ 的几何实现中, 这等价于:
\[
\phi(f) - \alpha_p f = 0
\]
在 $H^{1}\bigl(X_1(N),H_{k-2} \otimes \Omega^{\bullet}(\log)\bigr) \otimes L$ 中, 从而在 $H^{1}\bigl(W_1,H_{k-2} \otimes \Omega^{\bullet}(\log)\bigr) \otimes L$ 中 (因为限制是单射).
由引理 \ref{lem:4.3.2} 和对拟 Stein 开集 $W_1$ 应用的引理 \ref{lem:4.2.1}, 这也等价于:
\[
p^{k-1}\chi(p)f(pz) - \alpha_p f(z) = 0
\]
在 $H^{0}(W_1,\omega^{k})/\theta^{k-1}H^{0}(W_1,\omega^{2-k}) \otimes L$ 中, 最终即:
\[
f(z) - \beta_p f(pz) \in \theta^{k-1}H^{0}(W_1,\omega^{2-k}) \otimes L.
\]
现在假设 $1 \le N \le 4$.
如 \ref{sec:4.2}, 固定辅助素数 $q > 2$, $q \nmid Np$, 并用曲线 $X_1(N;q)$ 工作.
滤过 $\phi$-模 $D_p(f)$ 此时有几何描述:
\[
D_p(f) \simeq H^{1}\bigl(X_1(N;q),H_{k-2} \otimes \Omega^{\bullet}(\log)\bigr)^{f,\GL_2(\mathbb{Z}/q\mathbb{Z})} \otimes_{\Qp} L,
\]
\[
\mathrm{Fil}^{i}D_p(f) \simeq H^{0}(X_1(N;q),\omega^{k})^{f,\GL_2(\mathbb{Z}/q\mathbb{Z})} \otimes_{\Qp} L \quad (1 \le i \le k-1),
\]
即 $L \cdot f$ (因为 \cite{ref:29} 的比较定理对自同构 $\GL_2(\mathbb{Z}/q\mathbb{Z})$ 的作用是函子的).
用 $X_1(N;q)$ 代替 $X_1(N)$ 工作, 则如前证明 $\sigma_p(f)$ 分裂当且仅当存在 $g \in H^{0}(W_1,\omega^{2-k})^{\GL_2(\mathbb{Z}/q\mathbb{Z})} \otimes L$ 使得在 $H^{0}(W_1,\omega^{k})^{\GL_2(\mathbb{Z}/q\mathbb{Z})} \otimes L$ 中 $f(z) - \beta_p f(pz) = \theta^{k-1}(g)$.
\end{proof}

\subsection{(N,p) > 1 时定理 1.1.3 的证明}\label{sec:4.4}

保留 \ref{sec:4.3} 的记号与假设, 唯一区别是现在假设 $N = Mp^{r}$, $(M,p) = 1$, $r \ge 1$.
固定一个如 \ref{sec:3.3} 的新模形式 $f$, 使得 $\sigma_p(f)$ 在 $\Qp(\zeta_{p^{r}})$ 上变成晶体.
由于 $r > 0$, $f$ 是 $U_p$ 的特征向量, 再假设 $f$ 是零斜率的, 即相应特征值的 $p$ 进赋值是 $0$.
回顾在这种情况下 $\chi_1z^{k-2}$ 非分歧且 $U_p f = p^{k-1}\chi_1(p)f$ (参见 \ref{sec:4.1}).
若 $U$ 是定义在一个域上的簇 (代数或刚性) 且 $K$ 是该域的有限扩张, 记 $U_K$ 为到 $K$ 的基变换.
仍记 $H_{k-2} \otimes \Omega^{\bullet}(\log)$ 为 \ref{sec:4.2} 定义的复形到 $X_1(N)_K$ 的基变换.

设 $W_r \subset X_1(M)$ 为 $X_1(M)$ 的闭点 $x$ 满足 $|A(\bar{x})| > p^{-1/p^{r-2}(p+1)}$ 的刚性解析开集 (若 $r = 1$ 或 $r = 2$, 正好回到 $W_1$ 和 $W_2$, 参见 \ref{sec:4.3}).
若 $E$ 对应 $W_r$ 的一点 (忘记 $M$ 处级结构), 可以递推定义 $E^{(0)} = E$, 且对 $0 \le i \le r-1$, $E^{(i+1)} = E^{(i)}/\mathrm{can}$.
定义 $W_1(p^{r}) \subset X_1(Mp^{r})$ 为 $X_1(Mp^{r})$ 中对应三元组 $(E,Q_M,P)$ 的闭点的刚性开集 (其中 $Q_M$ (resp.\ $P$) 是 $E$ 上精确 $M$ 阶 (resp.\ $p^{r}$ 阶) 的有理点), 使得 $(E,Q_M)$ 对应 $X_1(M)$ 中属于 $W_r$ 的一个闭点, 且对 $0 \le i \le r-1$, $p^{r-1-i}P$ 在 $E^{(i)}$ 中的像生成 $E^{(i)}$ 的标准子群.

先假设 $M > 4$, 记 $\mathcal{X}_1(Mp^{r})$ 为 \cite{ref:36} 中构造的 $X_1(Mp^{r})_{\Qp(\zeta_{p^{r}})}$ 在 $\Zp[\zeta_{p^{r}}]$ 上表示模问题 $([\Gamma_1(M)],[\mathrm{bal.}\Gamma_1(p^{r})^{\mathrm{can}}])$ 的紧、平坦、正则模型.
由 \cite[\S13.12]{ref:36}, $\mathcal{X}_1(Mp^{r})$ 的特殊纤维是约化的, 是在 $\mathbb{F}_p$ 上 Igusa 曲线的超奇异点处互相相交的不交并.
恰有两个与 $\mathrm{Ig}(Mp^{r})$ 同构的不可约分支 (即按 \cite[\S12]{ref:36} 的意义表示 $\mathbb{F}_p$ 上模问题 $([\Gamma_1(M)],[\mathrm{Ig}(p^{r})])$ 的 Igusa 曲线).
容易验证其中一个, 记为 $\mathrm{Ig}_{\infty}$, 包含 $W_1(p^{r})_{\Qp(\zeta_{p^{r}})}$ 经特殊化箭头的像.
记 $\mathrm{Ig}_0$ 为另一个分支.

再假设 $M$ 足够大, 使 $X_1(Mp^{r})$ 和 $X_1(Mp^{r})$ 的特殊纤维的不可约分支的亏格都 $\ge 2$.
设 $K$ 为 $\Qp(\zeta_{p^{r}})$ 的有限扩张, $X_1(Mp^{r})_K$ 在 $K$ 上有稳定模型, $F_K$ 为其剩余域.
记 $X = X_1(Mp^{r})_K$, $\mathcal{X}$ 为 $X_1(Mp^{r})$ 从 $\Zp[\zeta_{p^{r}}]$ 到 $\mathcal{O}_K$ 的基变换, $\widetilde{\mathcal{X}}$ 为 $X_K$ 在 $\mathcal{O}_K$ 上的极小正则 (从而是稳定的) 模型.
由于 $\mathcal{X}$ 的特殊纤维中不可能有收缩, 有态射 $\widetilde{\mathcal{X}} \to \mathcal{X}$, 它是在特殊纤维上的一串爆破.
仍记 $\mathrm{Ig}_{\infty}$ (resp.\ $\mathrm{Ig}_0$) 为 $\widetilde{\mathcal{X}} \times_{\mathcal{O}_K} F_K$ 的双射地映到 $\mathrm{Ig}_{\infty} \times F_K$ (resp.\ $\mathrm{Ig}_0 \times F_K$) 的唯一不可约分支, 记 $W_{\infty} \subset X$ (resp.\ $W_0 \subset X$) 为 $\mathrm{Ig}_{\infty}$ (resp.\ $\mathrm{Ig}_0$) 经 $\widetilde{\mathcal{X}}$ 的特殊化箭头的逆像 (即管).
记 $\res_{\infty}$ (resp.\ $\res_0$) 为限制映射 $H^{1}\bigl(X,H_{k-2} \otimes \Omega^{\bullet}(\log)\bigr) \to H^{1}\bigl(W_{\infty},H_{k-2} \otimes \Omega^{\bullet}(\log)\bigr)$ (resp.\ 对 $W_0$).

\begin{lemma}\label{lem:4.4.1}
群 $H^{1}\bigl(W_{\infty},H_{k-2} \otimes \Omega^{\bullet}(\log)\bigr)$ 和 $H^{1}\bigl(W_0,H_{k-2} \otimes \Omega^{\bullet}(\log)\bigr)$ 自然配备 $\ell \nmid Mp$ 的 Hecke 算子 $T_{\ell}$ 和 $S_{\ell}$ 的作用, 它经 $\res_{\infty}$ 和 $\res_0$ 与 $H^{1}\bigl(X,H_{k-2} \otimes \Omega^{\bullet}(\log)\bigr)$ 上的作用交换.
\end{lemma}

\begin{proof}
由极小正则模型的唯一性, $X$ 的自同构 $\langle \ell \rangle$ ($\ell \nmid Mp$) 规范地延拓为 $\widetilde{\mathcal{X}}$ 及其特殊纤维的自同构, 后者保持 $\mathrm{Ig}_{\infty}$ 和 $\mathrm{Ig}_0$ (因为 $\mathcal{X}$ 的特殊纤维上就是如此).
于是 $X$ 的自同构 $\langle \ell \rangle$ 保持 $W_{\infty}$ 和 $W_0$, 由 $S_{\ell} = \ell^{k-2}\langle \ell \rangle$ 即得 $S_{\ell}$ 的断言.
记 $\mathcal{X}^{(\ell)}$ 为 \cite{ref:36} 中构造的 $\Zp[\zeta_{p^{r}}]$ 上表示模问题 $([\Gamma_1(M)],[\mathrm{bal.}\Gamma_1(p^{r})^{\mathrm{can}}],[\Gamma_0(\ell)])$ (即椭圆曲线 (广义) 之间带 $E$ 上 $([\Gamma_1(M)],[\mathrm{bal.}\Gamma_1(p^{r})^{\mathrm{can}}])$ 型级结构的次数 $\ell$ 同源 $E \to E'$) 的正则紧平坦概形.
设 $\mathrm{pr}_1$ (resp.\ $\mathrm{pr}_2$) 为投影 $\mathcal{X}^{(\ell)} \to \mathcal{X}$, 把 $(E \to E')$ 送到 $E$ (resp.\ 送到带诱导的 $([\Gamma_1(M)],[\mathrm{bal.}\Gamma_1(p^{r})^{\mathrm{can}}])$ 级结构的 $E'$).
如 \cite{ref:40} 那样证明两个纤维积 $\mathcal{X}^{(\ell)} \times_{\mathcal{X}} \widetilde{\mathcal{X}}$ (对 $\mathrm{pr}_1$ 和 $\mathrm{pr}_2$) 同构, 因为它们都同构于 $\mathcal{X}^{(\ell)}$ 在 $\mathcal{O}_K$ 上的极小正则 (且稳定) 模型.
记 $\widetilde{\mathrm{pr}}_1$ 和 $\widetilde{\mathrm{pr}}_2$ 为两个投影 $\mathcal{X}^{(\ell)} \times_{\mathcal{X}} \widetilde{\mathcal{X}} \to \widetilde{\mathcal{X}}$, 则 $\widetilde{\mathrm{pr}}_1^{-1}(\mathrm{Ig}_{\infty})$ (resp.\ $\widetilde{\mathrm{pr}}_1^{-1}(\mathrm{Ig}_0)$) 是 $\mathcal{X}^{(\ell)}$ 的稳定模型的特殊纤维的唯一不可约分支, 其在 $\mathcal{X}^{(\ell)}$ 中的像是投影到 $\widetilde{\mathcal{X}} \times_{\mathcal{O}_K} F_K$ 中 $\mathrm{Ig}_{\infty} \times F_K$ (resp.\ $\mathrm{Ig}_0 \times F_K$) 上的 Igusa 曲线 (经 $\mathrm{pr}_1$ 或 $\mathrm{pr}_2$).
因此 $\widetilde{\mathrm{pr}}_2\bigl(\widetilde{\mathrm{pr}}_1^{-1}(\mathrm{Ig}_{\infty})\bigr) = \mathrm{Ig}_{\infty}$ 于 $\widetilde{\mathcal{X}} \times_{\mathcal{O}_K} F_K$, 故在 $X$ 中 $\mathrm{pr}_2\bigl(\mathrm{pr}_1^{-1}(W_{\infty})\bigr) = W_{\infty}$ (resp.\ 对 $\mathrm{Ig}_0$ 和 $W_0$).
这正是按 $H^{1}\bigl(X,H_{k-2} \otimes \Omega^{\bullet}(\log)\bigr)$ 上定义 $T_{\ell}$ 的同样方式 (参见 \cite[\S8]{ref:12}) 由限制定义 $H^{1}\bigl(W_{\infty},H_{k-2} \otimes \Omega^{\bullet}(\log)\bigr)$ (resp.\ 对 $W_0$) 上自同态 $T_{\ell}$ 所需要的条件.
\end{proof}

由于 $X$ 在 $\mathcal{O}_K$ 上有半稳定约化, $K$-向量空间 $H^{1}\bigl(X,H_{k-2} \otimes \Omega^{\bullet}(\log)\bigr)$ 有一个规范 $K_0$-结构 $H^{1}_{\log\text{-}\mathrm{cris}}$ (其中 $K_0 \subseteq K$ 是最大的绝对非分歧子域), 配备仍记为 $\phi$ 的半线性 Frobenius 和满足 $N\phi = p\phi N$ 的单值算子 $N$.
由 \cite{ref:27} 和 \cite{ref:28} (另见 \cite{ref:16}), 有系数在 $L$ 中的滤过 $\phi$-模同构:
\begin{equation}
\label{eq:20}
\bigl(D_p(f)_{K_0},\ \phi,\ \mathrm{Fil}^{k-1}D_p(f)_K\bigr)
\simeq \bigl((H^{1}_{\log\text{-}\mathrm{cris}} \otimes_{\Qp} L)^{f},\ \phi \otimes 1,\ (H^{0}(X,\omega^{k}) \otimes_{\Qp} L)^{f}\bigr).
\end{equation}
设 $d := [K_0:\Qp]$, 记 $\phi_d$ 为 $H^{1}\bigl(X,H_{k-2} \otimes \Omega^{\bullet}(\log)\bigr)$ 上由 $H^{1}_{\log\text{-}\mathrm{cris}}$ 上的自同构 $\phi^{d}$ 从 $K_0$ 到 $K$ 的标量扩张定义的 $K$-线性自同构.
$K$-向量空间 $H^{1}\bigl(W_{\infty},H_{k-2} \otimes \Omega^{\bullet}(\log)\bigr)$ 和 $H^{1}\bigl(W_0,H_{k-2} \otimes \Omega^{\bullet}(\log)\bigr)$ 规范配备一个 $K$-线性 Frobenius $\Phi$, 与算子 $T_{\ell}$ 和 $S_{\ell}$ ($\ell \nmid Mp$) 交换, 且限制 $\res_{\infty}$ 和 $\res_0$ 与 $\Phi$ 和 $\phi_d$ 交换 (参见 \cite{ref:16}).

\begin{lemma}\label{lem:4.4.2}
限制 $\res_{\infty}$ 和 $\res_0$ 诱导同构:
\[
H^{1}\bigl(X,H_{k-2} \otimes \Omega^{\bullet}(\log)\bigr) \otimes L^{f}
\xrightarrow{\sim} H^{1}\bigl(W_{\infty},H_{k-2} \otimes \Omega^{\bullet}(\log)\bigr) \otimes L^{f}
\oplus H^{1}\bigl(W_0,H_{k-2} \otimes \Omega^{\bullet}(\log)\bigr) \otimes L^{f}.
\]
\end{lemma}

\begin{proof}
已知 $f$ 在 $p$ 处的特征标导子恰好是 $p^{r}$, 有包含:
\[
H^{1}\bigl(X,H_{k-2} \otimes \Omega^{\bullet}(\log)\bigr) \otimes L^{f} \subset H^{1}(X,\Sym^{k-2})^{\mathrm{prim}} \otimes L,
\]
其中 $H^{1}(X,\Sym^{k-2})^{\mathrm{prim}}$ 是 \cite[\S2]{ref:14} 中引入的群.
同构则容易由 \cite[Th.~2.1]{ref:14} 的同构:
\[
H^{1}(X,\Sym^{k-2})^{\mathrm{prim}} \otimes L \simeq H^{1}(W_{\infty},\Sym^{k-2})^{*} \otimes L \oplus H^{1}(W_0,\Sym^{k-2})^{*} \otimes L
\]
(它本身由 Mayer--Vietoris 正合列推出) 推出.
\end{proof}

设 $\zeta$ 为 $1$ 的原初 $p^{r}$ 次根.
回顾 $X$ 分类数据 $(E,Q_M,\pi,(P,Q))$, 其中 $E$ 是 $\mathcal{O}_K$-概形 $S$ 上的广义椭圆曲线, $Q_M$ 是 $E(S)$ 中精确 $M$ 阶的点 (按 \cite[\S1.4]{ref:36} 的意义), $\pi : E \to E'$ 是 $S$-同源且次数为 $p^{r}$, $P \in \mathrm{Ker}(\pi)(S)$ 是 $\mathrm{Ker}(\pi)$ 的 (按"Drinfel'd"意义的) 生成元, $Q \in \mathrm{Ker}(\pi^{*})(S)$ 是 $\mathrm{Ker}(\pi^{*})$ 的生成元 ($\pi^{*}$ 是 $\pi$ 的对偶同源), 满足 $\langle P,Q \rangle_{\pi} = \zeta$, 其中 $\langle\ \cdot,\cdot\ \rangle_{\pi}$ 是 $\mathrm{Ker}(\pi)$ 与 $\mathrm{Ker}(\pi^{*})$ 之间的规范交错配对 (参见 \cite[\S2.8]{ref:36}).
定义 $X$ 的自同构 $w_{\zeta}$ (它保持 $X$), 把 $(E,Q_M,\pi,(P,Q))$ 送到 $(E',\pi(Q_M),\pi^{*},(Q,-P))$.

\begin{lemma}\label{lem:4.4.3}
$X$ 的自同构 $w_{\zeta}$ 诱导同构 $w_{\zeta} : W_{\infty} \xrightarrow{\sim} W_0$ 和 $w_{\zeta} : W_0 \xrightarrow{\sim} W_{\infty}$.
\end{lemma}

\begin{proof}
由极小正则模型的唯一性, $w_{\zeta}$ 从 $X$ 规范地延拓到 $\widetilde{\mathcal{X}}$, 并容易验证它在 $X \times_{\mathcal{O}_K} F_K$ 中交换 $\mathrm{Ig}_{\infty}$ 和 $\mathrm{Ig}_0$, 由此得到结果.
\end{proof}

有交换图:
\[
\begin{CD}
E @>>> E\\
@VVV @VVV\\
X @>{w_{\zeta}}>> X
\end{CD}
\]
其中上面的水平箭头是复合:
\[
E \to E/\mathrm{Ker}(\pi) \simeq X \times_{w_{\zeta},X} E \to E,
\]
它在 $H^{1}\bigl(X,H_{k-2} \otimes \Omega^{\bullet}(\log)\bigr)$ 上诱导一个仍记为 $w_{\zeta}$ 的自同构.
由引理 \ref{lem:4.4.3}, 由此得到交换图:
\begin{equation}
\label{eq:21}
\begin{CD}
H^{1}\bigl(X,H_{k-2} \otimes \Omega^{\bullet}(\log)\bigr) @>{\res_0}>> H^{1}\bigl(W_0,H_{k-2} \otimes \Omega^{\bullet}(\log)\bigr)\\
@V{w_{\zeta}}VV @VV{w_{\zeta}}V\\
H^{1}\bigl(X,H_{k-2} \otimes \Omega^{\bullet}(\log)\bigr) @>{\res_{\infty}}>> H^{1}\bigl(W_{\infty},H_{k-2} \otimes \Omega^{\bullet}(\log)\bigr).
\end{CD}
\end{equation}

\begin{lemma}\label{lem:4.4.4}
设 $s \in H^{1}\bigl(X,H_{k-2} \otimes \Omega^{\bullet}(\log)\bigr)$ 满足 $\res_0(s) = 0$, 则 $\res_{\infty}(w_{\zeta}(s)) = 0$.
\end{lemma}

\begin{proof}
立即由 \eqref{eq:21} 得出.
\end{proof}

设 $\widetilde{W} := W_1(p^{r})_K \cap W_{\infty} \subset X$, 它是两个严格邻域的交, 因而是 $X$ 中序特殊化轨迹的一个严格邻域.
因此存在仿射型集 $W \subset \widetilde{W}$, 它仍是这个轨迹的严格邻域 (\cite[\S1.2]{ref:5}), 且有限制态射:
\begin{equation}
\label{eq:22}
H^{1}\bigl(W_{\infty},H_{k-2} \otimes \Omega^{\bullet}(\log)\bigr) \to H^{1}\bigl(W,H_{k-2} \otimes \Omega^{\bullet}(\log)\bigr).
\end{equation}

\begin{remark}\label{rem:4.4.5}
用与 \cite[\S8]{ref:13} 中相同的论证并使用 \cite{ref:2}, 可以证明限制 \eqref{eq:22} 实际上是同构.
我们不需要这个事实.
\end{remark}

记 $\res_{\infty}(f)$ (resp.\ $\res_0(f)$) 为直线 $K \otimes_{\Qp} L \cdot f \subset H^{1}\bigl(X,H_{k-2} \otimes \Omega^{\bullet}(\log)\bigr) \otimes L^{f}$ 经映射 $\res_{\infty}$ (resp.\ $\res_0$) 在 $H^{1}\bigl(W_{\infty},H_{k-2} \otimes \Omega^{\bullet}(\log)\bigr) \otimes L^{f}$ (resp.\ $H^{1}\bigl(W_0,H_{k-2} \otimes \Omega^{\bullet}(\log)\bigr) \otimes L^{f}$) 中的像.

\begin{proposition}\label{prop:4.4.6}
假设表示 $\sigma_p(f)$ 分裂, 则 $\res_0(f) = 0$.
\end{proposition}

\begin{proof}
恒有 $\res_{\infty}(f) \neq 0$, 否则经限制 \eqref{eq:22} 和对仿射型集 $W$ 应用的引理 \ref{lem:4.2.1}, $f$ 会在 $\theta^{k-1}$ (作用在权 $2-k$ 的超收敛形式空间上) 的像中, 这在 $f$ 于 $p$ 处零斜率时不可能 (像 \cite[Lem.~6.3]{ref:13} 那样进行).
假设 $\res_0(f)$ 非零.
则两个 $K \otimes_{\Qp} L$-自由模 $H^{1}\bigl(W_{\infty},H_{k-2} \otimes \Omega^{\bullet}(\log)\bigr) \otimes L^{f}$ 和 $H^{1}\bigl(W_0,H_{k-2} \otimes \Omega^{\bullet}(\log)\bigr) \otimes L^{f}$ 都非零.
由于 $H^{1}\bigl(X,H_{k-2} \otimes \Omega^{\bullet}(\log)\bigr) \otimes L^{f}$ 在 $K \otimes_{\Qp} L$ 上自由且秩为 $2$ (Eichler--Shimura), 由引理 \ref{lem:4.4.2} 它们必为秩 $1$.
由于它们被 $\phi_d$ 稳定, 由 \eqref{eq:20} 和 \ref{sec:3.3} 给出的 $D_p(f)$ 的描述 (必要时可能交换 $e_1$ 和 $e_2$):
\[
H^{1}\bigl(W_{\infty},H_{k-2} \otimes \Omega^{\bullet}(\log)\bigr) \otimes L^{f} \simeq K \otimes_{\Qp} L \cdot e_1,
\qquad
H^{1}\bigl(W_0,H_{k-2} \otimes \Omega^{\bullet}(\log)\bigr) \otimes L^{f} \simeq K \otimes_{\Qp} L \cdot e_2.
\]
由于 $\res_0(f) \neq 0$ 且 $\res_{\infty}(f) \neq 0$, 在 $\mathrm{Fil}^{k-1}D_p(f)_K \simeq K \otimes_{\Qp} L \cdot f$ 的描述 (参见 \ref{sec:3.3}) 中确实可见 $\delta \neq 0$, 这蕴含 $\sigma_p(f)$ 不分裂.
\end{proof}

回顾 $f\,|\,w_{p^{r}}$ 已在 \ref{sec:4.1} 中定义.

\begin{proposition}\label{prop:4.4.7}
若 $\sigma_p(f)$ 分裂, 则存在 $g \in H^{0}(W,\omega^{2-k}) \otimes_{\Qp} L$ 使得在 $H^{0}(W,\omega^{k}) \otimes_{\Qp} L$ 中 $f\,|\,w_{p^{r}} = \theta^{k-1}(g)$.
\end{proposition}

\begin{proof}
由命题 \ref{prop:4.4.6}, $\res_0(f) = 0$.
而已知 $w_{p^{r}}$ (差一个非零常数因子) 由 $H^{1}\bigl(X,H_{k-2} \otimes \Omega^{\bullet}(\log)\bigr)$ 上的自同构 $w_{\zeta}$ 诱导 (参见 e.g.\ 对 $r = 1$, $k = 2$ 情形的 \cite[\S6]{ref:33}).
由引理 \ref{lem:4.4.4}, 因此有 $\res_{\infty}(f\,|\,w_{p^{r}}) = 0$, 这由 \eqref{eq:22} 和对 $W$ 应用的引理 \ref{lem:4.2.1} 蕴含 $f\,|\,w_{p^{r}} = \theta^{k-1}(g)$.
\end{proof}

现在假设 $M$ 任意, 并如 \ref{sec:4.2} 固定一个辅助素数 $q > 2$, $q \nmid Mp$, 足够大使得 $X_1(M;q)$ 及其特殊纤维的不可约分支的亏格至少为 $2$.
有滤过 $\phi$-模同构:
\[
\bigl(D_p(f)_{K_0},\ \phi,\ \mathrm{Fil}^{k-1}D_p(f)_K\bigr)
\simeq \bigl((H^{1}_{\log\text{-}\mathrm{cris}} \otimes_{\Qp} L)^{f,\GL_2(\mathbb{Z}/q\mathbb{Z})},\ \phi \otimes 1,\ (H^{0}(X,\omega^{k}) \otimes_{\Qp} L)^{f,\GL_2(\mathbb{Z}/q\mathbb{Z})}\bigr).
\]
引理 \ref{lem:4.4.2} 和命题 \ref{prop:4.4.6} 通过把处处 $(\cdot)^{f}$ 换成 $(\cdot)^{f,\GL_2(\mathbb{Z}/q\mathbb{Z})}$ 很容易推广.
由此得到命题 \ref{prop:4.4.7} 的类似陈述: 当 $\sigma_p(f)$ 分裂时, 在 $H^{0}(W,\omega^{2-k})^{\GL_2(\mathbb{Z}/q\mathbb{Z})} \otimes L$ 中 $f\,|\,w_{p^{r}} = \theta^{k-1}(g)$.

\begin{theorem}\label{thm:4.4.8}
设 $f$ 是对 $\Gamma_1(Mp^{r})$ 为新的权 $k \ge 2$ 本征尖点形式, $(M,p) = 1$, $r > 0$.
假设 $f$ 在 $p$ 处零斜率.
若 $Mp^{r} > 4$, 则表示 $\sigma_p(f)$ 分裂当且仅当存在 $g \in H^{0}(W_1(p^{r}),\omega^{2-k}) \otimes_{\Qp} L$ 使得在 $H^{0}(W_1(p^{r}),\omega^{k}) \otimes_{\Qp} L$ 中 $f\,|\,w_{p^{r}} = \theta^{k-1}(g)$.
若 $Mp^{r} \le 4$, 取 $q \nmid Mp$ 为足够大的素数, 则 (用 $X_1(N;q)$ 代替 $X_1(N)$) $\sigma_p(f)$ 分裂当且仅当存在 $g \in H^{0}(W_1(p^{r}),\omega^{2-k})^{\GL_2(\mathbb{Z}/q\mathbb{Z})} \otimes_{\Qp} L$ 使得在 $H^{0}(W_1(p^{r}),\omega^{k})^{\GL_2(\mathbb{Z}/q\mathbb{Z})} \otimes_{\Qp} L$ 中 $f\,|\,w_{p^{r}} = \theta^{k-1}(g)$.
\end{theorem}

\begin{proof}
若存在如陈述中的 $g$, 则已知 $\sigma_p(f)$ 分裂 (参见 e.g.\ \cite[Prop.~11]{ref:30}).
假设 $\sigma_p(f)$ 分裂, 并先假设不需要辅助素数 (即所有亏格至少为 $2$).
则已由命题 \ref{prop:4.4.7} 得到 $g \in H^{0}(W,\omega^{2-k}) \otimes_{\Qp} L$ 使得 $f\,|\,w_{p^{r}} = \theta^{k-1}(g)$.
只需证明 $g$ (必唯一地) 延拓为 $H^{0}(W_1(p^{r})_K,\omega^{2-k}) \otimes_{\Qp} L$ 的截面 (它实际上会在 $H^{0}(W_1(p^{r}),\omega^{2-k}) \otimes_{\Qp} L$ 中).
若 $k > 2$, 则 $\theta^{k-1}$ 是单射, 关系 $p\theta \circ U_p = U_p \circ \theta$ (在超收敛形式空间中) 蕴含 $U_p g = \lambda g$ 且 $\val(\lambda) = 0$, 因为 $f\,|\,w_{p^{r}}$ 斜率为 $k-1$.
连续注入 $H^{0}(W_1(p^{r})_K,\omega^{2-k}) \otimes L \hookrightarrow H^{0}(W,\omega^{2-k}) \otimes L$ 的像稠密 (注意左边的空间是 Fr\'echet 空间, 右边是 Banach 空间).
取 $(g_n)_n \in H^{0}(W_1(p^{r})_K,\omega^{2-k}) \otimes L$ 使得在 $H^{0}(W,\omega^{2-k}) \otimes L$ 中 $g_n \to g$, 记 $e := \lim U_p^{n!}$ 为定义在超收敛形式空间上的 Hida 投影子.
$e\bigl(H^{0}(W_1(p^{r})_K,\omega^{2-k}) \otimes L\bigr)$ 在 $H^{0}(W_1(p^{r})_K,\omega^{2-k}) \otimes L$ 中闭, 因为在 $K$ 上有限维 (因为 $U_p$ 是 $H^{0}(W_1(p^{r})_K,\omega^{2-k})$ 上的紧算子, 参见 \cite{ref:14}).
而 $e(g_n) \to e(g) = g$, 故 $g \in H^{0}(W_1(p^{r})_K,\omega^{2-k}) \otimes L$.
若 $k = 2$, 我们只能先验地有 $U_p g = \lambda g + \text{常数}$, 但同样的论证适用, 因为仍有 $e(g) = g$.

在我们引入辅助素数 $q$ 的情形, 用曲线 $X_1(Mp^{r};q)$ 工作, 由前述讨论同样的证明成立.
最后假设 $Mp^{r} > 4$, 使得 $X_1(Mp^{r})$ 是精细模空间, 为避免混淆, 记 $W_1(p^{r};q)$ 为 $X_1(Mp^{r};q)$ 中此前一直简记 $W_1(p^{r})$ 的刚性解析开集.
这个子空间是 $\GL_2(\mathbb{Z}/q\mathbb{Z})$-不变的, 其经 $\GL_2(\mathbb{Z}/q\mathbb{Z})$ 作用的商同构于 $X_1(Mp^{r})$ 的开集 $W_1(p^{r})$.
此外有自然同构:
\[
H^{0}(W_1(p^{r};q),\omega^{2-k})^{\GL_2(\mathbb{Z}/q\mathbb{Z})} \otimes_{\Qp} L \simeq H^{0}(W_1(p^{r}),\omega^{2-k}) \otimes_{\Qp} L,
\]
\[
H^{0}(W_1(p^{r};q),\omega^{k})^{\GL_2(\mathbb{Z}/q\mathbb{Z})} \otimes_{\Qp} L \simeq H^{0}(W_1(p^{r}),\omega^{k}) \otimes_{\Qp} L
\]
(参见 \cite{ref:10} 的附录).
在这种情形下, 可以如陈述那样删去对辅助素数 $q$ 的一切提及.
\end{proof}
